% META: {"id": "TSPE-LOG-ME-002", "chapitre": "TSPE-LOGARITHME", "type_objet": "methode", "capacites_codes": ["C3", "C4"], "methodes": ["M2"], "sources_inspiration": [], "status": "approved"}
\begin{fichemethode}{M2}{Étudier une fonction construite avec $\ln$}

\quandUtiliser{%
  Utiliser cette méthode pour dériver, étudier les limites, ou dresser le tableau de variations d'une fonction faisant intervenir $\ln$.
}

\pasApas{%
  \begin{enumerate}
    \item \textbf{Domaine} : déterminer l'ensemble de définition (condition de positivité sur l'argument du $\ln$).
    \item \textbf{Dérivée} : utiliser $(\ln u)' = \dfrac{u'}{u}$ pour une composee, ou les règles usuelles pour une somme/produit.
    \item \textbf{Limites aux bornes} : utiliser $\ln x \to +\infty$ en $+\infty$, $\ln x \to -\infty$ en $0^+$, et au besoin les croissances comparees ($x\ln x \to 0$ en $0^+$, $\ln x / x \to 0$ en $+\infty$).
    \item \textbf{Signe de la dérivée et tableau de variations.}
  \end{enumerate}
}

\end{fichemethode}
